% GENERATED from content/data/enrichment.json by mxbuild.enrichment; do not edit.
\mxunit{enrichment-bisection-solutions}{Encadrer une racine par dichotomie}
\mxsubtitle{Terminale · analyse et algorithmique \textperiodcentered{} Solutions détaillées \textperiodcentered{} 12}

\begin{note}
\textbf{Données.} Pour l’exercice 1 : on cherche un zéro de $f(x)=x^{2}-2$ sur $\intff{1}{2}$. Après deux partages, l’intervalle conservé est $\intff{\num{1,25}}{\num{1,5}}$.
\end{note}

\begin{exercise}
Poursuivre l’exemple une troisième fois.
\end{exercise}
\begin{solution}
Le milieu vaut $\num{1,375}$ ; son carré vaut $\num{1,890625}<2$. Le nouvel intervalle est $\intff{\num{1,375}}{\num{1,5}}$.
\end{solution}

\begin{exercise}
Partant d’un intervalle de largeur 1, combien de partages garantissent une largeur au plus $\frac{1}{1000}$ ?
\end{exercise}
\begin{solution}
$2^{9}=512<1000\leq 1024=2^{10}$, donc 10 partages suffisent et 9 ne suffisent pas.
\end{solution}

\begin{exercise}
Pourquoi le changement de signe de $f(x)=\frac{1}{x}$ entre $-1$ et 1 ne prouve-t-il pas l’existence d’un zéro ?
\end{exercise}
\begin{solution}
La fonction n’est pas définie en 0 et n’est pas continue sur $\intff{-1}{1}$. L’hypothèse du théorème n’est pas satisfaite ; $\frac{1}{x}$ n’est jamais nul.
\end{solution}

\begin{exercise}
Une racine est dans $\intff{\num{1,4}}{\num{1,5}}$. Quelle erreur maximale garantit le milieu $\num{1,45}$ ?
\end{exercise}
\begin{solution}
La demi-largeur vaut $\frac{\num{1,5}-\num{1,4}}{2}=\num{0,05}$.
\end{solution}

\begin{problem}[Chercher, justifier, conclure]
Pour $f(x)=x^{3}-x-1$ sur $\intff{1}{2}$, prouver l’existence et l’unicité d’une racine $\alpha $. Effectuer trois partages par dichotomie. Combien de partages garantissent que le milieu final approche $\alpha $ avec une erreur au plus $\num{0,001}$ ?
\end{problem}
\begin{solution}
$f(1)=-1$ et $f(2)=5$ ; f est continue. De plus $f'(x)=3x^{2}-1\geq 2$ sur $\intff{1}{2}$, donc f est strictement croissante. Les milieux $\num{1,5}$, $\num{1,25}$, $\num{1,375}$ donnent successivement $\intff{1}{\num{1,5}}$, $\intff{\num{1,25}}{\num{1,5}}$, $\intff{\num{1,25}}{\num{1,375}}$ car les valeurs de f sont $\num{0,875}$, $-\num{0,296875}$ et $\num{0,224609375}$. Après n partages, l’erreur du milieu est au plus $\frac{1}{2^{n+1}}$. Il faut $2^{n+1}\geq 1000$, donc $n\geq 9$.
\end{solution}

